\documentclass[11pt]{article}
\usepackage[a4paper,margin=1in]{geometry}
\usepackage{fontspec}
\usepackage[boldfont=false]{xeCJK}
\setCJKmainfont{FandolSong-Regular.otf}
\usepackage{amsmath}
\usepackage{amssymb}
\usepackage{array}
\usepackage{booktabs}
\usepackage{graphicx}
\usepackage{adjustbox}
\usepackage{enumitem}
\setlist[itemize]{topsep=2pt,partopsep=0pt,parsep=0pt,itemsep=2pt,leftmargin=1.4em}
\setlist[enumerate]{topsep=2pt,partopsep=0pt,parsep=0pt,itemsep=2pt,leftmargin=1.6em}
\usepackage{parskip}
\usepackage{fvextra}
\fvset{breaklines=true,breakanywhere=true,fontsize=\small}
\setlength{\parindent}{0pt}

\begin{document}

\begin{center}
  {\LARGE\bfseries Operations management}\\[0.35em]
  {\large Igcse Business Studies}
\end{center}
\vspace{0.6em}

\section*{Methods of production}

\textbf{Production} 生产 means turning materials and \textbf{labour} 劳动力 into finished goods and \textbf{services} 服务. There are three main methods.

\begin{center}
\adjustbox{max width=\linewidth}{%
\begin{tabular}{llll}
\toprule
\textbf{Method} & \textbf{What it is} & \textbf{Good for} & \textbf{Drawback} \\
\midrule
\textbf{job production} 单件生产 & making one item at a time, to order & special, one-off items: a wedding cake, a ship & slow and costly per item \\
\textbf{batch production} 批量生产 & making a group of the same item, then switching to another group & bakeries, clothes in different sizes & machines stand idle while switching \\
\textbf{flow production} 流水线生产 & making goods non-stop on a moving line & mass-made goods: cars, drinks & costly to set up; work can be boring \\
\bottomrule
\end{tabular}}
\end{center}

\subsection*{Productivity}

\textbf{Productivity} 生产率 measures how much each worker or machine produces in a set time. Higher productivity lowers the cost per item. A business can raise productivity by:

\begin{itemize}
\item training workers,

\item using better machines and \textbf{technology} 技术,

\item improving \textbf{motivation} 激励,

\item organising the work better.

\end{itemize}

\subsection*{Lean production and just-in-time}

\textbf{Lean production} 精益生产 means cutting all waste — wasted time, materials, space and effort — while keeping quality.

One key method is \textbf{just-in-time} 准时制 (JIT). With JIT, materials arrive just as they are needed, and finished goods are sold quickly, so the business holds very little \textbf{inventory} 库存 (its stock of materials and goods). This saves storage cost, but one late delivery can stop production. Good management of inventory keeps enough stock to meet demand, but not so much that money is tied up or goods go out of date.

\subsection*{Technology in production}

Technology has changed production. \textbf{Automation} 自动化 (machines and robots doing tasks), computer design and online ordering all make production faster, cheaper and more exact. But machines cost a lot to buy and can replace workers' jobs.

\section*{Costs, scale and break-even}

\subsection*{Types of cost}

A business has different kinds of \textbf{cost} 成本.

\begin{itemize}
\item \textbf{fixed costs} 固定成本 do not change with output — rent, salaries, insurance. You pay them even if you make nothing.

\item \textbf{variable costs} 可变成本 change with \textbf{output} 产量 — materials, and pay for each item made. More output means more variable cost.

\item \textbf{total cost} 总成本 is fixed costs plus variable costs.

\item \textbf{average cost} 平均成本 is the total cost divided by the number of items made.

\end{itemize}

\[
\text{total cost} = \text{fixed costs} + \text{variable costs}
\]

\[
\text{average cost} = \frac{\text{total cost}}{\text{output}}
\]

\subsection*{Revenue and profit}

\textbf{Revenue} 收入 is the money a business gets from selling its products.

\[
\text{total revenue} = \text{price} \times \text{quantity sold}
\]

\textbf{Profit} 利润 is what is left after all costs are paid.

\[
\text{profit} = \text{total revenue} - \text{total cost}
\]

\subsection*{Break-even}

The \textbf{break-even} 盈亏平衡 point is the level of output where total revenue exactly equals total cost: the business makes no profit and no loss. Below it the business makes a loss; above it, a profit.

Each item sold gives an amount towards covering the fixed costs. This amount is the selling price minus the variable cost of one item (sometimes called the \textbf{contribution} 贡献 per unit).

\[
\text{break-even output} = \frac{\text{fixed costs}}{\text{selling price per unit} - \text{variable cost per unit}}
\]

A \textbf{break-even chart} 盈亏平衡图 plots total cost and total revenue against output. The two lines cross at the \textbf{break-even point} 盈亏平衡点.

The \textbf{margin of safety} 安全边际 is how far the actual output is above the break-even output.

\[
\text{margin of safety} = \text{actual output} - \text{break-even output}
\]

A large margin of safety is safer, because sales can fall a long way before the business makes a loss.

\textbf{Limits of break-even analysis.} It assumes that all output is sold, and that the selling price and the costs stay the same at every level of output. In real life these often change, so break-even is only a guide.

\section*{Achieving quality}

\subsection*{Why quality matters}

\textbf{Quality} 质量 means a product is good enough for what \textbf{customers} 顾客 need. Good quality:

\begin{itemize}
\item builds a strong \textbf{brand} 品牌 and \textbf{customer loyalty} 顾客忠诚度,

\item means fewer returns, complaints and wasted materials,

\item lets a business charge a higher price.

\end{itemize}

Poor quality loses customers and damages the brand.

\subsection*{Quality control and quality assurance}

There are two main ways to manage quality: \textbf{quality control} 质量控制 and \textbf{quality assurance} 质量保证.

\begin{center}
\adjustbox{max width=\linewidth}{%
\begin{tabular}{lll}
\toprule
\textbf{} & \textbf{Quality control} & \textbf{Quality assurance} \\
\midrule
When & check at the end & build in quality at every stage \\
Who & inspectors check finished goods & every worker is responsible \\
Idea & find and remove faulty goods & stop faults from happening at all \\
\bottomrule
\end{tabular}}
\end{center}

\section*{Location decisions}

\subsection*{Locating a business}

Where a business sits affects its costs and its sales. A \textbf{manufacturing} 制造业 business (one that makes goods) looks at:

\begin{itemize}
\item nearness to \textbf{raw materials} 原材料, because heavy or bulky materials cost a lot to move,

\item nearness to customers and the \textbf{market} 市场,

\item the cost and size of the land,

\item a supply of suitable labour,

\item good transport links,

\item government grants offered in some areas.

\end{itemize}

A \textbf{service} business (one that sells services, such as a shop or a hair salon) looks more at being close to customers, passing foot traffic, nearby shops, and the rent.

\subsection*{Locating in another country}

A business may produce or sell abroad in order to reach new customers and bigger markets, pay lower wages or rent, get nearer to raw materials, or avoid \textbf{tariffs} 关税 (taxes on imports) by making goods inside the country. But it must then deal with different laws, languages and customs, and follow the \textbf{legal controls} 法律管制 of both countries.


\end{document}
